\section{Final Project Report Draft}\label{final-project-report-draft}

Generated: 2026-05-08

This report summarizes the current project state from existing artifacts
only. No new RL training was run while preparing this document.

\subsection{Executive Summary}\label{executive-summary}

This project studies whether market-regime information can improve
weekly portfolio allocation across \texttt{SPY}, \texttt{TLT},
\texttt{GLD}, and cash. The original pipeline used an HMM regime model,
but the active research path now uses a PCA Jump Model because it gives
direct control over regime-switching frequency, supports an
interpretable K sweep, and produces compact state features for RL and
attention models.

The current leak-safe Jump Model pipeline builds weekly price and macro
features, applies causal rolling robust standardization, compresses the
feature set with train-fitted PCA, fits Jump Model centroids on the
train split only, and assigns validation/test regimes with online causal
rules. A 12-week sequence dataset is then exported for attention or RL
training.

Main findings:

\begin{itemize}
\tightlist
\item
  The selected smoothed PCA Jump Model uses 6 PCs, K=3, jump penalty
  32.0, and explains 53.94\% of feature variance. The compression ratio
  is 54 original numeric market/macro features to 6 PCs, or 9.0:1.
\item
  The selected regimes are economically interpretable: calm/risk-on,
  growth/trend, and stress/risk-off. The stress regime has the highest
  VIX and deepest SPY drawdown profile.
\item
  K=2 has the best raw silhouette, while K=5 is the elbow choice. K=3 is
  the recommended reporting choice because it balances interpretability,
  duration, minimum cluster size, and attention-readiness.
\item
  Mutual information supports using volatility, drawdown, rates, and
  bond-market features as regime/RL inputs, but the signal is modest and
  target dependent.
\item
  Existing RL results are mixed. A quick Jump RL run performs well on
  validation and modestly on locked test, but the long DQN tuning
  artifact underperforms simple baselines. This is an important result,
  not just a failure: regime detection can be useful descriptively while
  the trading policy still needs stronger action design, reward design,
  and validation discipline.
\end{itemize}

\subsection{Project Evolution}\label{project-evolution}

The proposal began with HMM-style market-regime detection and RL
allocation. During experimentation, several HMM-specific assumptions
became limiting:

\begin{itemize}
\tightlist
\item
  Gaussian emission assumptions made the states sensitive to scaling and
  outliers.
\item
  Prior filtering rules constrained the model before the data structure
  was fully understood.
\item
  Regime persistence was indirect, controlled through transition
  probabilities rather than a clear switch penalty.
\item
  Full-sample standardization risked making early years look
  artificially quiet compared with later high-variance or high-volume
  periods.
\end{itemize}

The Jump Model path addresses these issues by clustering in PCA space
while penalizing label changes through an explicit jump penalty. This
makes it easier to tune the model around market-regime behavior: not
only cluster separation, but also switch frequency, average duration,
and usefulness for the downstream attention/RL state.

\subsection{Data And Splits}\label{data-and-splits}

The core state table is:

\begin{itemize}
\tightlist
\item
  \texttt{data/processed/model\_state\_weekly\_price\_macro.csv}
\end{itemize}

The train-ready Jump Model artifacts are:

\begin{itemize}
\tightlist
\item
  \texttt{data/processed/jump\_model\_train\_ready\_weekly.csv}
\item
  \texttt{data/processed/jump\_model\_train\_ready\_sequences.csv}
\item
  \texttt{data/processed/jump\_model\_train\_ready\_sequences.npz}
\item
  \texttt{data/processed/jump\_model\_train\_ready\_metadata.json}
\end{itemize}

The split design is chronological:

{\def\LTcaptype{none} % do not increment counter
\begin{longtable}[]{@{}lrll@{}}
\toprule\noalign{}
split & rows & start\_week & end\_week \\
\midrule\noalign{}
\endhead
\bottomrule\noalign{}
\endlastfoot
train & 406 & 2014-03-28 & 2021-12-31 \\
validation & 104 & 2022-01-07 & 2023-12-29 \\
locked test & 115 & 2024-01-05 & 2026-03-13 \\
\end{longtable}
}

The sequence dataset uses a 12-week lookback. A validation or test
sequence may look back into prior history, which is acceptable for
online inference because that history would already be known. The split
assignment is based on the sequence end week, and no target values from
the future are included in features.

\subsection{Leakage-Safe
Preprocessing}\label{leakage-safe-preprocessing}

Leakage control is one of the most important report topics for this
project. The active train-ready dataset avoids the major leakage paths
as follows:

\begin{itemize}
\tightlist
\item
  Target columns such as \texttt{next\_return\_spy},
  \texttt{next\_return\_tlt}, \texttt{next\_return\_gld}, and
  \texttt{best\_asset\_next\_week} are excluded from \texttt{x\_*}
  features.
\item
  PCA is fit on the train split only.
\item
  Jump Model centroids are fit on the train split only.
\item
  Validation and locked-test regimes are assigned using current and past
  information only.
\item
  Continuous RL features are standardized using train-split statistics
  after the causal Jump Model features are built.
\item
  Soft-score temperature and regime naming/VIX ordering are derived from
  train assignments only.
\end{itemize}

The rolling raw-feature scaler is causal, not full-sample. For week
\texttt{t}, the rolling robust scaler uses trailing observed history up
to that point, with a 52-week window, at least 12 weeks of history, and
clipping at +/-6.0. This is appropriate for online inference because a
deployed system can always look back at past market history. It is
different from full-sample standardization, which would leak future
distribution shifts into earlier observations.

\subsection{Why Rolling Robust Scaling Was
Needed}\label{why-rolling-robust-scaling-was-needed}

The first Jump Model visualizations showed many early short-lived
regimes and later regimes dominated by higher variance and larger
changes in volume or macro levels. This is consistent with a common
clustering failure mode: nonstationary feature scale can dominate
distance-based models.

The rolling robust scaler reduces this issue by standardizing each
feature relative to its recent local history. Median/MAD-style scaling
is less sensitive to spikes than mean/std scaling, and the clip value
prevents extreme values from overpowering PCA and cluster distances. The
goal is not to erase market stress; the goal is to let the model
distinguish unusual moves relative to the market context available at
that time.

\subsection{PCA Compression}\label{pca-compression}

The selected Jump Model uses 6 PCA components from 54 numeric
market/macro features:

{\def\LTcaptype{none} % do not increment counter
\begin{longtable}[]{@{}lr@{}}
\toprule\noalign{}
item & value \\
\midrule\noalign{}
\endhead
\bottomrule\noalign{}
\endlastfoot
original features & 54 \\
PCA components & 6 \\
compression ratio & 9.0:1 \\
explained variance & 53.94\% \\
\end{longtable}
}

The compression ratio is a mechanical ratio: original feature count
divided by PCA component count. It should not be optimized alone. In
this project, 6 PCs are preferred because they preserve more structure
for downstream attention than 2 PCs, while still reducing noise and
collinearity enough for stable clustering.

\subsection{Jump Model Selection}\label{jump-model-selection}

The selected research Jump Model is:

{\def\LTcaptype{none} % do not increment counter
\begin{longtable}[]{@{}lr@{}}
\toprule\noalign{}
parameter & value \\
\midrule\noalign{}
\endhead
\bottomrule\noalign{}
\endlastfoot
scaler & rolling robust \\
scaler window & 52 weeks \\
scaler min history & 12 weeks \\
scaler clip & +/-6.0 \\
PCA components & 6 \\
K & 3 \\
jump penalty & 32.0 \\
post-cluster smoothing & 3-week minimum run \\
\end{longtable}
}

Selected-K metrics:

{\def\LTcaptype{none} % do not increment counter
\begin{longtable}[]{@{}lr@{}}
\toprule\noalign{}
metric & value \\
\midrule\noalign{}
\endhead
\bottomrule\noalign{}
\endlastfoot
inertia & 26732.84 \\
silhouette & 0.1844 \\
jumps & 29 \\
min duration & 3 weeks \\
average duration & 19.30 weeks \\
max duration & 95 weeks \\
\end{longtable}
}

The K sweep shows why this should not be chosen by silhouette alone:

{\def\LTcaptype{none} % do not increment counter
\begin{longtable}[]{@{}rrrrr@{}}
\toprule\noalign{}
K & inertia & silhouette & jumps & average duration \\
\midrule\noalign{}
\endhead
\bottomrule\noalign{}
\endlastfoot
2 & 31094.35 & 0.3256 & 24 & 23.16 \\
3 & 26732.84 & 0.1844 & 29 & 19.30 \\
4 & 23927.78 & 0.1670 & 27 & 20.68 \\
5 & 21622.99 & 0.1728 & 26 & 21.44 \\
6 & 19851.10 & 0.1761 & 34 & 16.54 \\
7 & 18346.26 & 0.1643 & 39 & 14.48 \\
8 & 17105.95 & 0.1734 & 42 & 13.47 \\
9 & 16622.98 & 0.1686 & 40 & 14.12 \\
10 & 15920.72 & 0.1666 & 43 & 13.16 \\
\end{longtable}
}

K=2 wins raw silhouette because it separates the sample into broad
low-stress and high-stress areas. K=5 is the elbow-selected value. K=3
is more useful for attention/RL because it keeps a distinct stress state
while avoiding too many thin regimes.

\subsection{Regime Interpretation}\label{regime-interpretation}

Regime IDs are ordered from lowest average VIX to highest average VIX.

{\def\LTcaptype{none} % do not increment counter
\begin{longtable}[]{@{}
  >{\raggedright\arraybackslash}p{(\linewidth - 12\tabcolsep) * \real{0.1154}}
  >{\raggedright\arraybackslash}p{(\linewidth - 12\tabcolsep) * \real{0.1154}}
  >{\raggedleft\arraybackslash}p{(\linewidth - 12\tabcolsep) * \real{0.1538}}
  >{\raggedleft\arraybackslash}p{(\linewidth - 12\tabcolsep) * \real{0.1538}}
  >{\raggedleft\arraybackslash}p{(\linewidth - 12\tabcolsep) * \real{0.1538}}
  >{\raggedleft\arraybackslash}p{(\linewidth - 12\tabcolsep) * \real{0.1538}}
  >{\raggedleft\arraybackslash}p{(\linewidth - 12\tabcolsep) * \real{0.1538}}@{}}
\toprule\noalign{}
\begin{minipage}[b]{\linewidth}\raggedright
regime
\end{minipage} & \begin{minipage}[b]{\linewidth}\raggedright
interpretation
\end{minipage} & \begin{minipage}[b]{\linewidth}\raggedleft
weeks
\end{minipage} & \begin{minipage}[b]{\linewidth}\raggedleft
share
\end{minipage} & \begin{minipage}[b]{\linewidth}\raggedleft
VIX
\end{minipage} & \begin{minipage}[b]{\linewidth}\raggedleft
SPY 20d return
\end{minipage} & \begin{minipage}[b]{\linewidth}\raggedleft
next SPY annualized
\end{minipage} \\
\midrule\noalign{}
\endhead
\bottomrule\noalign{}
\endlastfoot
R0 & Calm / risk-on & 309 & 53.37\% & 15.44 & 1.93\% & 16.33\% \\
R1 & Growth / trend & 172 & 29.71\% & 17.46 & 2.26\% & 21.02\% \\
R2 & Stress / risk-off & 98 & 16.93\% & 25.29 & -4.37\% & -6.49\% \\
\end{longtable}
}

This broadly aligns with real-world market behavior. The stress/risk-off
regime has higher VIX, deeper SPY drawdown, higher SPY volatility, and
negative forward SPY returns. The calm and growth regimes both support
risk assets, but the growth state shows stronger trend behavior and
higher forward SPY return.

The report should be careful not to claim perfect event detection. These
are statistical states derived from price/macro features, not named
macro events. The stronger claim is that the states have economically
coherent profiles.

\subsection{Smoothing And Duration}\label{smoothing-and-duration}

There are two smoothing concepts in the project:

\begin{enumerate}
\def\labelenumi{\arabic{enumi}.}
\tightlist
\item
  Descriptive post-clustering smoothing for charts and static Jump Model
  artifacts. Runs shorter than the configured minimum duration are
  merged into the closest adjacent regime by PCA-centroid distance. This
  reduces 1-2 week visual noise in regime timelines.
\item
  Causal confirmation smoothing for RL features. A new regime must
  persist for the configured number of weeks before the confirmed regime
  switches. This creates a delay, but it does not inspect future weeks
  and is therefore acceptable for online inference.
\end{enumerate}

This distinction should be highlighted in the final report. Static
smoothing is for interpretation; causal smoothing is for train-ready RL
state construction.

\subsection{Mutual Information
Evidence}\label{mutual-information-evidence}

Mutual information was used to check whether market/macro features
contain signal for next-week allocation targets. The analysis used 579
rows and 54 numeric features, excluding close levels and future returns.

For the primary classification target, \texttt{best\_asset\_next\_week},
the top features were:

{\def\LTcaptype{none} % do not increment counter
\begin{longtable}[]{@{}rlrr@{}}
\toprule\noalign{}
rank & feature & MI & permutation p-value \\
\midrule\noalign{}
\endhead
\bottomrule\noalign{}
\endlastfoot
1 & \texttt{tnx\_level} & 0.0407 & 0.0198 \\
2 & \texttt{gld\_vol\_20d} & 0.0338 & 0.0594 \\
3 & \texttt{dgs10\_level} & 0.0289 & 0.0495 \\
4 & \texttt{tlt\_ret\_1d} & 0.0281 & 0.0990 \\
5 & \texttt{spy\_drawdown\_60d} & 0.0267 & 0.0495 \\
\end{longtable}
}

For SPY next return, the strongest signals were
\texttt{spy\_drawdown\_60d}, \texttt{vix\_level},
\texttt{spy\_intraday\_range}, \texttt{umcsent\_level}, and
\texttt{spy\_vol\_20d}. For TLT, bond volatility and the yield curve
were more important. For GLD, TLT returns and rate features were more
informative.

Interpretation: the signal is real but modest. This supports using these
features in a regime-aware state representation, but it also explains
why a high-capacity RL agent can overfit if the action and reward setup
is not tightly controlled.

\subsection{RL-Ready Dataset}\label{rl-ready-dataset}

The train-ready dataset is designed for attention and RL:

\begin{itemize}
\tightlist
\item
  12-week lookback sequences.
\item
  21 leak-safe features.
\item
  PCA features, regime scores, centroid-distance information, regime
  duration, regime-change indicators, and target-only next-week returns.
\item
  Targets: \texttt{y\_next\_return\_spy}, \texttt{y\_next\_return\_tlt},
  \texttt{y\_next\_return\_gld}, \texttt{y\_best\_asset\_id}, and
  \texttt{y\_best\_asset}.
\end{itemize}

The attention context is important: the model should not only see the
current regime label. It should also see the trajectory into that
regime, distances to alternative regimes, and whether a regime
transition is newly forming or already persistent.

\subsection{RL Setup}\label{rl-setup}

The Jump RL environment uses weekly allocation actions over:

{\def\LTcaptype{none} % do not increment counter
\begin{longtable}[]{@{}ll@{}}
\toprule\noalign{}
action & allocation \\
\midrule\noalign{}
\endhead
\bottomrule\noalign{}
\endlastfoot
\texttt{cash\_only} & 100\% cash \\
\texttt{spy\_only} & 100\% SPY \\
\texttt{tlt\_only} & 100\% TLT \\
\texttt{gld\_only} & 100\% GLD \\
\texttt{spy\_80\_tlt\_20} & 80\% SPY / 20\% TLT \\
\texttt{balanced\_60\_30\_10} & 60\% SPY / 30\% TLT / 10\% GLD \\
\texttt{defensive\_20\_60\_20} & 20\% SPY / 60\% TLT / 20\% GLD \\
\end{longtable}
}

The evaluation protocol selects models using validation metrics only.
Locked test is reported after selection and must not be used to choose
hyperparameters.

\subsection{Existing RL Results}\label{existing-rl-results}

Two relevant RL result sets exist in the repository.

\subsubsection{Quick Jump RL Runner}\label{quick-jump-rl-runner}

Artifact:

\begin{itemize}
\tightlist
\item
  \texttt{output/full\_pipeline/jump\_rl\_training\_summary.json}
\end{itemize}

Configuration:

\begin{verbatim}
python scripts/train_jump_rl.py \
  --timesteps 30000 --seed 123 --min-action-hold-weeks 8 \
  --reward-scale 100.0 --risk-penalty 0.05
\end{verbatim}

Metrics are net of transaction costs:

{\def\LTcaptype{none} % do not increment counter
\begin{longtable}[]{@{}
  >{\raggedright\arraybackslash}p{(\linewidth - 12\tabcolsep) * \real{0.1111}}
  >{\raggedleft\arraybackslash}p{(\linewidth - 12\tabcolsep) * \real{0.1481}}
  >{\raggedleft\arraybackslash}p{(\linewidth - 12\tabcolsep) * \real{0.1481}}
  >{\raggedleft\arraybackslash}p{(\linewidth - 12\tabcolsep) * \real{0.1481}}
  >{\raggedleft\arraybackslash}p{(\linewidth - 12\tabcolsep) * \real{0.1481}}
  >{\raggedleft\arraybackslash}p{(\linewidth - 12\tabcolsep) * \real{0.1481}}
  >{\raggedleft\arraybackslash}p{(\linewidth - 12\tabcolsep) * \real{0.1481}}@{}}
\toprule\noalign{}
\begin{minipage}[b]{\linewidth}\raggedright
split
\end{minipage} & \begin{minipage}[b]{\linewidth}\raggedleft
cumulative return
\end{minipage} & \begin{minipage}[b]{\linewidth}\raggedleft
annualized return
\end{minipage} & \begin{minipage}[b]{\linewidth}\raggedleft
volatility
\end{minipage} & \begin{minipage}[b]{\linewidth}\raggedleft
Sharpe
\end{minipage} & \begin{minipage}[b]{\linewidth}\raggedleft
max drawdown
\end{minipage} & \begin{minipage}[b]{\linewidth}\raggedleft
turnover
\end{minipage} \\
\midrule\noalign{}
\endhead
\bottomrule\noalign{}
\endlastfoot
validation & 19.56\% & 9.34\% & 15.49\% & 0.603 & -18.58\% & 0.115 \\
locked test & 9.65\% & 4.25\% & 14.37\% & 0.296 & -11.90\% & 0.113 \\
\end{longtable}
}

This run is promising because it improves validation performance while
keeping turnover low. The locked-test result is positive but much
weaker, so it should be reported as preliminary rather than definitive.
Because this is a limited runner artifact rather than the full long-DQN
tuning sweep, it should be framed as an exploratory result.

\subsubsection{Long DQN Tuning Artifact}\label{long-dqn-tuning-artifact}

Artifact:

\begin{itemize}
\tightlist
\item
  \texttt{output/jump\_rl\_long\_dqn/best\_config.json}
\end{itemize}

Best selected trial:

\begin{itemize}
\tightlist
\item
  algorithm: DQN
\item
  seed: 7
\item
  timesteps: 2,000,000
\item
  learning rate: 0.0001
\item
  buffer size: 250,000
\item
  batch size: 256
\item
  gamma: 0.99
\end{itemize}

Metrics:

{\def\LTcaptype{none} % do not increment counter
\begin{longtable}[]{@{}
  >{\raggedright\arraybackslash}p{(\linewidth - 10\tabcolsep) * \real{0.1304}}
  >{\raggedleft\arraybackslash}p{(\linewidth - 10\tabcolsep) * \real{0.1739}}
  >{\raggedleft\arraybackslash}p{(\linewidth - 10\tabcolsep) * \real{0.1739}}
  >{\raggedleft\arraybackslash}p{(\linewidth - 10\tabcolsep) * \real{0.1739}}
  >{\raggedleft\arraybackslash}p{(\linewidth - 10\tabcolsep) * \real{0.1739}}
  >{\raggedleft\arraybackslash}p{(\linewidth - 10\tabcolsep) * \real{0.1739}}@{}}
\toprule\noalign{}
\begin{minipage}[b]{\linewidth}\raggedright
split
\end{minipage} & \begin{minipage}[b]{\linewidth}\raggedleft
cumulative return
\end{minipage} & \begin{minipage}[b]{\linewidth}\raggedleft
annualized excess return
\end{minipage} & \begin{minipage}[b]{\linewidth}\raggedleft
Sharpe
\end{minipage} & \begin{minipage}[b]{\linewidth}\raggedleft
max drawdown
\end{minipage} & \begin{minipage}[b]{\linewidth}\raggedleft
turnover
\end{minipage} \\
\midrule\noalign{}
\endhead
\bottomrule\noalign{}
\endlastfoot
validation & -7.85\% & -6.37\% & -0.430 & -23.91\% & 0.500 \\
locked test & -5.97\% & -6.31\% & -0.438 & -22.33\% & 0.570 \\
\end{longtable}
}

This long run underperforms simple baselines and has high turnover. That
result suggests the current DQN setup is not yet robust, even if the
regime model itself is meaningful.

\subsection{Baselines}\label{baselines}

The strongest locked-test baselines are difficult to beat:

{\def\LTcaptype{none} % do not increment counter
\begin{longtable}[]{@{}llrrr@{}}
\toprule\noalign{}
strategy & split & cumulative return & Sharpe & max drawdown \\
\midrule\noalign{}
\endhead
\bottomrule\noalign{}
\endlastfoot
\texttt{momentum\_rotation\_20d} & validation & 16.88\% & 0.394 &
-11.64\% \\
\texttt{gld\_only} & validation & 12.76\% & 0.258 & -17.35\% \\
\texttt{gld\_only} & locked test & 118.10\% & 1.770 & -14.55\% \\
\texttt{equal\_weight\_spy\_tlt\_gld} & locked test & 47.14\% & 1.439 &
-8.32\% \\
\texttt{momentum\_rotation\_20d} & locked test & 60.61\% & 1.080 &
-14.86\% \\
\end{longtable}
}

The 2024-2026 locked-test period strongly favors GLD, which makes static
GLD and simple baseline comparisons unusually strong. The report should
therefore avoid claiming that the RL model is generally superior. A
better conclusion is that regime-aware state construction is promising,
but policy learning remains the main bottleneck.

\subsection{Streamlit Dashboard}\label{streamlit-dashboard}

The Streamlit dashboard supports real-time-style market replay and
regime diagnostics:

\begin{itemize}
\tightlist
\item
  market replay with indexed SPY/TLT/GLD prices
\item
  cluster scatter view
\item
  PCA component visualization
\item
  regime time-series timeline
\item
  elbow/K sweep diagnostics
\item
  regime summary table
\item
  heatmap-style regime feature profiles
\end{itemize}

Research defaults reflected in the dashboard:

{\def\LTcaptype{none} % do not increment counter
\begin{longtable}[]{@{}lr@{}}
\toprule\noalign{}
control & default \\
\midrule\noalign{}
\endhead
\bottomrule\noalign{}
\endlastfoot
scaler window weeks & 52 \\
scaler minimum history weeks & 12 \\
scaler clip & 6.0 \\
jump penalty & 6.0 \\
minimum displayed regime duration & 6 \\
K sweep & 2 to 10 \\
K selection & manual \\
manual K & 4 \\
\end{longtable}
}

The dashboard is useful for research exploration, but final RL artifacts
should be produced by the leak-safe pipeline scripts, not by full-sample
interactive refits.

\subsection{Topics That Should Not Be Missing From The Final
Report}\label{topics-that-should-not-be-missing-from-the-final-report}

The final submitted report should explicitly include these topics:

{\def\LTcaptype{none} % do not increment counter
\begin{longtable}[]{@{}
  >{\raggedright\arraybackslash}p{(\linewidth - 2\tabcolsep) * \real{0.5000}}
  >{\raggedright\arraybackslash}p{(\linewidth - 2\tabcolsep) * \real{0.5000}}@{}}
\toprule\noalign{}
\begin{minipage}[b]{\linewidth}\raggedright
topic
\end{minipage} & \begin{minipage}[b]{\linewidth}\raggedright
why it matters
\end{minipage} \\
\midrule\noalign{}
\endhead
\bottomrule\noalign{}
\endlastfoot
HMM-to-Jump-Model pivot & Explains why the project changed method rather
than silently switching. \\
Leakage controls & Prevents the most common critique of financial ML
projects. \\
Rolling robust scaling & Addresses nonstationary variance/volume scale
and the user's observed early-regime noise. \\
PCA compression ratio & Shows the theoretical feature compression and
why 6 PCs were selected. \\
K selection tradeoff & Explains why raw silhouette was not the only
objective. \\
Jump penalty and duration & Shows regimes are intended to be persistent
market states, not weekly noise. \\
Static vs causal smoothing & Clarifies what is safe for RL and what is
only descriptive. \\
Mutual information & Provides feature-target evidence before RL. \\
Baseline comparison & Keeps the RL results honest. \\
Locked-test discipline & Shows test was not used to choose model
settings. \\
Failure analysis & Turns weak long-DQN performance into a research
insight. \\
Streamlit dashboard & Demonstrates inspection, replay, and
interpretability. \\
Future RL improvements & Gives a credible path beyond current
results. \\
\end{longtable}
}

\subsection{Limitations}\label{limitations}

\begin{itemize}
\tightlist
\item
  Regime labels are statistical and should not be treated as
  ground-truth macro states.
\item
  Silhouette is low for K\textgreater2 because market regimes overlap;
  low silhouette does not automatically invalidate the regimes, but it
  limits confidence.
\item
  The locked-test period is unusual because GLD performs extremely well.
\item
  The long DQN artifact underperforms baselines, suggesting overfitting,
  poor exploration, action-template mismatch, or reward instability.
\item
  Transaction costs and turnover matter materially; high-frequency
  regime switching can destroy performance.
\item
  The current asset universe is small. Adding more ETFs may improve
  allocation opportunities but also increases action complexity.
\end{itemize}

\subsection{Future Work}\label{future-work}

Recommended next steps:

\begin{enumerate}
\def\labelenumi{\arabic{enumi}.}
\tightlist
\item
  Keep the leak-safe Jump Model dataset as the canonical regime feature
  source.
\item
  Compare RL against a simple supervised policy that predicts the best
  next-week asset from the same leak-safe features.
\item
  Add regime-aware but rule-based allocation baselines, such as stress
  regime to cash/GLD and calm/growth regime to SPY.
\item
  Tune action templates, especially reducing all-in actions and adding
  smoother allocation choices.
\item
  Train attention models with explicit regime score and regime duration
  inputs, then inspect attention weights by market period.
\item
  Evaluate across rolling walk-forward windows so the conclusion is not
  too dependent on the 2024-2026 GLD-dominated test period.
\end{enumerate}

\subsection{Reproducibility}\label{reproducibility}

Build the train-ready Jump Model dataset:

\begin{verbatim}
python scripts/build_train_ready_dataset.py
\end{verbatim}

Run the Streamlit dashboard:

\begin{verbatim}
streamlit run app/streamlit_jump_model.py
\end{verbatim}

Quick Jump RL command represented by the existing summary artifact:

\begin{verbatim}
python scripts/train_jump_rl.py \
  --timesteps 30000 --seed 123 --min-action-hold-weeks 8 \
  --reward-scale 100.0 --risk-penalty 0.05
\end{verbatim}

Long DQN tuning artifacts are already stored under:

\begin{verbatim}
output/jump_rl_long_dqn/
\end{verbatim}

Key supporting reports:

\begin{itemize}
\tightlist
\item
  \texttt{reports/jump\_model\_train\_ready\_dataset.md}
\item
  \texttt{reports/jump\_model\_tuning\_rolling\_robust52\_smoothed.md}
\item
  \texttt{reports/mutual\_information\_results.md}
\item
  \texttt{reports/jump\_rl\_tuning.md}
\end{itemize}
